\documentclass[11pt,a4paper]{article}

\usepackage[utf8]{inputenc}
\usepackage[T1]{fontenc}
\usepackage[brazil]{babel}
\usepackage{lmodern}
\usepackage{geometry}
\usepackage{microtype}
\usepackage{xcolor}
\usepackage{tabularx}
\usepackage{booktabs}
\usepackage{enumitem}
\usepackage{tikz}
\usepackage{fancyhdr}
\usepackage{hyperref}

\usetikzlibrary{arrows.meta,positioning}
\geometry{margin=2.3cm, headheight=14pt}
\setlength{\parindent}{0pt}
\setlength{\parskip}{0.6em}
\setlist[itemize]{leftmargin=1.5em,itemsep=0.2em}

\newcommand{\jogo}{The Last Contract}
\definecolor{Destaque}{HTML}{7A2E2E}

\pagestyle{fancy}
\fancyhf{}
\fancyhead[L]{\textcolor{Destaque}{\jogo}}
\fancyhead[R]{Projeto}
\fancyfoot[C]{\thepage}
\renewcommand{\headrulewidth}{0.4pt}
\hypersetup{colorlinks=true, linkcolor=black, urlcolor=Destaque, pdftitle={\jogo{} - Projeto}}

\title{\textbf{\jogo} \\ \large Projeto de jogo}
\author{André Bartzen Araujo \and Jackson Matiazzi \and José Maciel Camargo}
\date{Programação de Jogos --- Avaliação Grau 2}

\begin{document}

\maketitle
\thispagestyle{empty}

\section{Conceito}

Jogo de ação 3D em terceira pessoa ambientado em uma vila medieval tomada por mortos-vivos. O jogador escolhe um entre quatro aventureiros e aceita o contrato de um Lorde: eliminar os mortos-vivos e derrotar o responsável pela infestação.

Cada tentativa tem uma vida só. A proposta é explorar a vila livremente, enfrentar grupos de inimigos e decidir, a cada derrota, se vale a pena continuar.

\section{História}

Uma vila medieval foi completamente dominada por mortos-vivos. Diante da ameaça, um Lorde convoca aventureiros para eliminar a infestação e derrotar a criatura que a provocou: um necromante escondido na vila.

O contrato tem uma cláusula dura. Se o aventureiro morrer, pode aceitar uma nova tentativa. Se desistir, o Lorde considera o abandono uma falha grave e ordena sua execução em praça pública.

\section{Gameplay}

O jogador explora a vila e enfrenta grupos de mortos-vivos espalhados pelo mapa. A cada tentativa, o ponto de partida do jogador e as posições dos grupos são sorteados, então o caminho e os encontros mudam.

Ao morrer, o jogador escolhe entre \textbf{reviver} ou \textbf{desistir}. Ao reviver, todos os inimigos voltam e o jogador escolhe uma entre duas bênçãos aleatórias. A bênção vale para o aventureiro e também para os mortos-vivos nas tentativas seguintes, o que transforma cada escolha numa aposta.

\begin{center}
\begin{tikzpicture}[
    node distance=0.9cm and 1.2cm,
    etapa/.style={draw=Destaque, rounded corners, thick, minimum width=3.4cm,
                  minimum height=0.9cm, align=center, font=\small},
    fim/.style={etapa, fill=Destaque!12},
    seta/.style={-{Stealth[length=2.5mm]}, thick, Destaque},
    rotulo/.style={font=\footnotesize, midway}
]
    \node[etapa] (contrato) {Aceitar o contrato\\do Lorde};
    \node[etapa, below=of contrato] (sorteio) {Sortear posições do\\jogador e dos inimigos};
    \node[etapa, below=of sorteio] (combate) {Explorar a vila e\\enfrentar os grupos};
    \node[etapa, below=of combate] (boss) {Derrotar o boss};
    \node[fim, below=of boss] (vitoria) {Vitória};

    \node[etapa, right=2.2cm of combate] (morte) {Morte};
    \node[etapa, above=of morte] (buff) {Escolher 1 entre\\2 bênçãos};
    \node[fim, below=of morte] (execucao) {Desistir:\\execução};

    \draw[seta] (contrato) -- (sorteio);
    \draw[seta] (sorteio) -- (combate);
    \draw[seta] (combate) -- (boss);
    \draw[seta] (boss) -- (vitoria);
    \draw[seta] (combate) -- (morte);
    \draw[seta] (boss.east) -- (morte.south west);
    \draw[seta] (morte) -- node[rotulo, right] {Reviver} (buff);
    \draw[seta] (buff) |- (sorteio);
    \draw[seta] (morte) -- (execucao);
\end{tikzpicture}
\end{center}

\section{Mecânicas}

\begin{tabularx}{\textwidth}{@{}lX@{}}
\toprule
\textbf{Mecânica} & \textbf{Descrição} \\
\midrule
Movimentação & Terceira pessoa, com andar, correr e pular. A câmera segue o personagem por trás do ombro. \\
Combate & Corpo a corpo ou à distância, conforme o aventureiro. Cada um tem um ataque principal e uma habilidade própria. \\
Vida única & Uma vida por tentativa. Ao morrer, reviver ou desistir. \\
Bênçãos & Ao reviver, o jogador escolhe uma entre duas bênçãos; elas se acumulam e valem também para os inimigos. \\
Grupos de inimigos & Mortos-vivos organizados em grupos. Ao derrotar um grupo inteiro, os restos podem se juntar num inimigo mais forte. \\
Posições aleatórias & Jogador e grupos nascem em pontos sorteados a cada tentativa. \\
Objetos interativos & Baús espalhados pela vila guardam itens de cura. \\
Boss & O necromante surge depois que os grupos da vila caem. \\
Vitória & Derrotar todos os grupos e o boss. \\
\bottomrule
\end{tabularx}

\section{Personagens}

\textbf{Aventureiros jogáveis.} Cada um tem um estilo de combate:

\begin{tabularx}{\textwidth}{@{}llX@{}}
\toprule
\textbf{Aventureiro} & \textbf{Alcance} & \textbf{Estilo} \\
\midrule
Knight & Corpo a corpo & Espada e escudo; combo de golpes e defesa com o escudo. \\
Barbarian & Corpo a corpo & Dois machados; giro que acerta todos em volta. \\
Mage & Longo alcance & Projétil mágico que atravessa inimigos em linha. \\
Rogue & Longo alcance & Besta precisa, de longo alcance. \\
\bottomrule
\end{tabularx}

\textbf{Lorde.} Contrata o aventureiro e julga o abandono do contrato.

\textbf{Mortos-vivos.} Esqueletos que patrulham, dormem no chão e despertam com o jogador por perto, perseguem quem veem ou ouvem e procuram por ele onde foi visto pela última vez.

\textbf{Necromante.} O boss: um mago cercado de mortos-vivos, responsável pela infestação.

\section{Mundo}

Toda a missão acontece num único mapa: uma vila medieval com ruas, casas, praça e um campo aberto, explorada à noite. Não há fases separadas; a variação entre tentativas vem do sorteio das posições do jogador e dos grupos de inimigos.

\end{document}
